\section{Two models}\label{app:models}

Section~\ref{sec:structure} reads two measured quantities through two textbook models.

\paragraph{Ornstein--Uhlenbeck process}~\cite{uhlenbeck1930} (Section~\ref{sec:cohesion}). A quantity
$\ell$ pulled back to 0 by a spring of rate $k$ and shaken by noise:
\begin{equation}
  d\ell=-k\,\ell\,dt+\sqrt{2k}\,\sigma\,dW .
\end{equation}
It forgets its value after $1/k$ and spreads by $\sigma$ around 0. Its mean square change after a
lag $s$ is $D(s)=\langle(\ell(t+s)-\ell(t))^2\rangle=2\sigma^2(1-e^{-ks})$: it grows as
$2k\sigma^2 s$ at short lags, like a free random walk, and stops at $2\sigma^2$ after $1/k$. With
$k=0$ it is a random walk and $D$ grows forever. A measurement noise of variance $c/2$ on each
value adds $c$ to $D(s)$ at every lag.

\paragraph{Birth--death chain} (Section~\ref{sec:fission}). A group of $n$ pilots gains one at rate
$g(n)$ and loses one at rate $l(n)$. Over a long time the transitions from $n$ to $n+1$ and back are
equally many, so the times spent at the two sizes satisfy $T(n)\,g(n)=T(n+1)\,l(n+1)$,
Eq.~\eqref{eq:chain}, and multiplying the ratios from 2 to $n$ gives $T(n)$. If every member brings
in others at rate $a$ and leaves at rate $\mu$, then $g(n)=an$, $l(n)=\mu n$ and $T(n)\propto
x^n/n$ with $x=a/\mu$, the log-series, whose tail falls exponentially. A leaving rate that falls
with $n$ makes the tail heavier.
